\documentclass[11pt]{article}

\usepackage[margin=1in]{geometry}
\usepackage{amsmath,amssymb,amsthm,mathtools}
\usepackage{enumitem}
\usepackage{hyperref}

\hypersetup{
  colorlinks=true,
  linkcolor=blue,
  citecolor=blue,
  urlcolor=blue
}

\newtheorem{definition}{Definition}
\newtheorem{proposition}{Proposition}
\newtheorem{theorem}{Proposition}
\newtheorem{assumption}{Assumption}
\theoremstyle{remark}
\newtheorem{remark}{Remark}

\newcommand{\Effects}{\mathcal{E}}
\newcommand{\LTL}{\mathrm{LTL}}
\newcommand{\LTLf}{\mathrm{LTL}_f}
\newcommand{\AP}{\mathrm{AP}}
\newcommand{\Close}{\mathrm{Close}}
\newcommand{\Lang}{\mathcal{L}}
\newcommand{\Nat}{\mathbb{N}}
\newcommand{\True}{\mathsf{true}}
\newcommand{\Frame}{\mathfrak{F}}
\newcommand{\Design}{\eta}
\newcommand{\Adm}{\mathrm{Adm}}

\title{Frame-Relative Counterfactual Necessity for Explaining Synthesized
Temporal Behavior}
\author{}
\date{}

\begin{document}
\maketitle

\begin{abstract}
LTL synthesis is many-to-one: a controller alone fixes neither a unique
specification nor a unique temporal explanation of its behavior. We study the
reverse question under a declared \emph{causal frame}: given background,
intervention model, candidate vocabulary, similarity order, and preference order,
when is a candidate \(C\) a \emph{frame-relative counterfactual necessary
condition} for effect \(E\)?

We separate two notions throughout. \textbf{Trace-existential necessity (N1)}
applies to a Mealy machine \(M\), trace \(\tau\), and frame: deciding validity
for a fixed \(C\) is a one-player shortest-edit search plus non-emptiness check
(the frozen machine has no free move). N1 is confined to co-safety / finite-change
causes; \(\omega\)-liveness commitments require the strategy level. \textbf{Strategy-guarantee
necessity (N2)} additionally requires a synthesis artifact \((G_\varphi,\sigma,V_{\Frame})\)
(not determined by \(M\) alone): a constrained \emph{reverse parity game} decides
whether the controller can drop \(C\) and still guarantee \(E\) against admissible
environments (parity winning region, then a one-player check that re-imposes the
objective).

For a \emph{fixed} candidate \(C\) and finite candidate space
\(\mathcal{K}_{\Frame}\), validity and \(\triangleleft\)-minimality are decidable
under stated bounds; we do \emph{not} present a complete solver or a tight
complexity theorem. Canonical default frames fix every component except the
declared vocabulary, grammar, and tie-breaks; the answer is then a function of
\((M,\tau,E)\) \emph{given that declaration}. A finite-trace specialization
aligns with TempoBench-style TCE \emph{conditionally} on machine-checkable
assumptions: a control-validated audit confirms them on the publicly available
TempoBench sample artifact, and a reference implementation of the finite
construction reproduces five worked instances from HOA encodings and matches
the one available TempoBench answer key per cell. Corpus-scale evaluation
awaits the dataset's re-release.
\end{abstract}

\tableofcontents

\section{Introduction}
\label{sec:intro}

The standard synthesis pipeline is
\[
  \varphi_{\LTL}\longrightarrow A_{\varphi}\longrightarrow G_{\varphi}
  \longrightarrow \sigma \longrightarrow M ,
\]
where \(\varphi_{\LTL}\) is an LTL specification, \(A_{\varphi}\) an automaton for
its models, \(G_{\varphi}\) a parity game, \(\sigma\) a winning strategy, and
\(M=(S,s_0,I,O,\delta,\lambda)\) a Mealy implementation (\(\delta:S\times
2^I\to S\), \(\lambda:S\times 2^I\to 2^O\)). The pipeline loses information: many
specifications yield the same machine and many strategies are behaviorally
equivalent, so the reconstructive reverse map \(M\to\varphi\) is ill-posed.

We pose instead an \emph{explanatory} reverse question. The inputs split by
regime. \textbf{Trace-level analysis} takes \((M,\Frame,E,\tau)\): a Mealy
machine \(M=(S,s_0,I,O,\delta,\lambda)\) fixing the I/O partition (\(I\cap
O=\emptyset\)), a causal frame \(\Frame\), a temporal effect \(E\) (a single,
possibly conjunctive, \(\omega\)-regular trace property), and observed trace
\(\tau\in\Lang(M)\). \textbf{Strategy-level analysis} additionally requires a
synthesis/game artifact \((G_\varphi,\sigma,V_{\Frame})\) --- parity game
\(G_\varphi=(V,V_0,V_1,T,\Omega)\), frozen winning strategy \(\sigma\) implemented
by \(M\), and editable vertex set \(V_{\Frame}\subseteq V_0\) --- which \(M\) alone
does not determine. Optional specification \(\varphi\) may justify \(G_\varphi\) and
\(\Frame=\operatorname{Frame}(\varphi;\Design)\), but the strategy query is not a
function of \((M,\tau,E)\) alone. Throughout we use ``cause'' for a
\emph{frame-relative necessary condition}; this is counterfactual necessity, not
Halpern--Pearl actual causation (Section~\ref{sec:overdetermination}).

The task is to \emph{define and check} candidates and, when
\(\mathcal{K}_{\Frame}\) is finite or bounded, to \emph{search} for a
\(\triangleleft\)-minimal valid one \(C^\star\); we do not claim a complete solver
for open LTL candidate spaces or minimal dominion extraction. A trace is the infinite word
\(\tau=a_0a_1\cdots\) with \(a_t=i_t\cup o_t\in 2^{I\cup O}\),
\(o_t=\lambda(s_t,i_t)\), \(s_{t+1}=\delta(s_t,i_t)\); in implementations \(\tau\)
is a lasso \(uv^\omega\).

The intended output is not ``the controller wins,'' but a temporal explanation of
what had to remain true for the chosen effect to occur, \emph{relative to a
declared frame}. We commit to a \emph{counterfactual-necessity} reading: \(C\)
holds on \(\tau\), removing \(C\) is realizable under the frame, and on the
closest admissible executions that remove \(C\) the effect fails. This is not
sufficiency and not Halpern--Pearl actual causation (which adds a
minimal-sufficient-set requirement); Section~\ref{sec:overdetermination} discusses
what the necessity reading excludes. We avoid the unqualified phrase ``weakest
necessary condition'': in standard logic the \emph{strongest} necessary condition
for \(E\) (equivalently the weakest \emph{sufficient} condition) is \(E\) itself,
while the weakest necessary condition is \(\True\), so the phrase is doubly
misleading for what we do. Instead, ``more permissive'' means larger language
\(\Lang(C)\), ``insufficient'' means ``does not by itself force the effect,'' and
these are kept distinct.

\section{Related Work and Contributions}
\label{sec:related}

This work sits at the intersection of three lines of research; our claims are
deltas relative to them.

\paragraph{Counterfactual and actual causation.} The necessity core is
Lewis--Stalnaker: a cause is what, on the closest worlds where it is absent,
removes the effect, with centering and a limit assumption on the similarity
order. Halpern--Pearl actual causation adds a minimal-sufficient-set requirement
and a contingency quantifier; we adopt only the necessity half
(Section~\ref{sec:overdetermination}) and are explicit about what this excludes
(preemption discrimination, sufficiency).

\paragraph{Temporal/\(\omega\)-regular causality.} The closest technical
relatives are automata-theoretic accounts of causes for \(\omega\)-regular
effects in reactive systems (downward-closed \(\omega\)-regular causes via
game/automaton constructions), which also underlie TempoBench ground truth. Our
delta is not ``another parity solve'' but \emph{what the game is asked}:
\begin{itemize}[nosep]
  \item \textbf{Centering:} N1 is trace-existential (closest admissible
        counterfactual trace); N2 is strategy-centered (can \(\sigma\) drop a
        commitment while still guaranteeing \(E\)?), without Lewis--Stalnaker
        trace centering.
  \item \textbf{Objective:} \(B_{\Frame}\rightarrow E\) with non-commitment
        witnessed on \(B_{\Frame}\wedge\neg C\) plays, not a single global
        \(\omega\)-regular cause language.
  \item \textbf{Intervention:} re-opening a declared vertex subset
        \(V_{\Frame}\) of a \emph{frozen} \(\sigma\), not synthesizing an
        unconstrained alternative controller.
  \item \textbf{Framing:} background, admissible counterfactuals, similarity,
        candidate vocabulary, and anti-circularity are explicit parameters; prior
        constructions fix many of these implicitly.
\end{itemize}
We do not claim a new causality-synthesis algorithm; the contribution is this
parameterization plus sound decision procedures for fixed \(C\) (and bounded
search when \(|\mathcal{K}_{\Frame}|\) is finite).

\paragraph{Reactive synthesis and parity games.} We use only standard facts:
LTL\(\to\)deterministic parity automaton translation; parity games are determined,
solvable in quasi-polynomial time, and lie in \(\mathrm{UP}\cap\mathrm{co\text{-}UP}\);
HOA is the interchange format for \(\omega\)-automata. These are background.

\paragraph{Contributions.} (1)~An explicit \emph{causal frame} \(\Frame\)
factoring a temporal causal query into background, effect scope, intervention
model, candidate vocabulary, similarity order, and preference order, with
canonical default frames (Section~\ref{sec:canonical-frames}) that pin all
components but the declared vocabulary, grammar, and tie-breaks. (2)~A clean
separation of trace-existential necessity (N1: one-player shortest-edit +
non-emptiness) from strategy-guarantee necessity (N2: reverse parity game with
sound-and-complete decision for fixed \(C\), Section~\ref{sec:rpg}). (3)~A
decidable anti-circularity condition (Assumption~\ref{ass:anticirc}) --- semantic
position-masking for point/bounded effects, pure-logic non-entailment for all
effects, and a mandatory effect cone for general \(\omega\)-regular effects ---
that bars logical entailment and point-position restatement while admitting
input/invariant causes. (4)~A finite-trace specialization aligning with
TempoBench-style TCE \emph{conditionally} on (F1)--(F2), which we have not
verified (Section~\ref{sec:tempobench}). Core deliverables are the definitions
and Propositions~\ref{prop:rpg-sound}--\ref{prop:tb-reduction}; a tight
complexity theorem, dominion minimization, and empirical evaluation are left open.

\paragraph{Scope (read once).} \textbf{Established:} valid-cause clauses; N1/N2
distinction; Proposition~\ref{prop:rpg-sound} (N2 decision for fixed
\(V_{\Frame},C\)); certificate characterization
(Proposition~\ref{prop:weakest}); finite specialization
(Proposition~\ref{prop:tb-reduction}). \textbf{Scoped:} candidate checking and
search over finite/bounded \(\mathcal{K}_{\Frame}\); conditional complexity
(Section~\ref{sec:complexity}); TempoBench alignment on (F1)--(F2).
\textbf{Not claimed:} full solver, HP-style actual causation, verified benchmark
artifacts, or tight minimal-dominion completeness.

\section{Background: Parity Games}
\label{sec:background}

A parity game is \(G=(V,V_0,V_1,T,\Omega)\) with vertices \(V=V_0\sqcup V_1\)
(system \(V_0\), environment \(V_1\)), a total edge relation \(T\subseteq V\times
V\), and a priority map \(\Omega:V\to\Nat\). A play is an infinite path; under the
\emph{min-even} convention used throughout, the system wins iff the least priority
seen infinitely often is even. LTL formulas translate to deterministic parity
automata, so synthesis reduces to deciding \(\exists\sigma\,\forall\rho.\
\mathrm{Play}(\sigma,\rho)\models\varphi\); a winning \(\sigma\) is implementable
as a Mealy machine. We reuse this machinery in reverse: we freeze \(M\) as the
strategy and ask which temporal conditions were necessary for a chosen effect.

\section{The Formal Framework}
\label{sec:formal}

\subsection{The causal frame and admissible counterfactuals}

The reverse-causal object is the frame
\[
  \Frame=(B_{\Frame},E_{\Frame},S_{\Frame},O_{\Frame},
          \mathcal{K}_{\Frame},\mathcal{I}_{\Frame},{\preceq},\triangleleft),
\]
where \(B_{\Frame}\) is a background constraint (environment assumptions and
system guarantees, when present, are folded in; \(B_{\Frame}=\True\) if none),
\(E_{\Frame}=\bigwedge_{r\in S_{\Frame}}E_r\) is the effect selected from a
decomposition \(E=E_1\land\cdots\land E_m\) by indices \(S_{\Frame}\) (if no
scoping, \(E_{\Frame}=E\)), \(O_{\Frame}\subseteq O\) the relevant outputs,
\(\mathcal{K}_{\Frame}\) the candidate-cause space, \(\mathcal{I}_{\Frame}\) the
intervention model, \({\preceq}=\{\preceq_\tau\}_\tau\) the similarity orders, and
\(\triangleleft\) the candidate-preference order (\(C'\triangleleft C\): \(C'\) is
preferred). When derived from a specification we write
\(\Frame=\operatorname{Frame}(\varphi;\Design)\); \(\Frame\) is \emph{not}
\(\varphi\) and is unrelated to a Kripke frame. A useful slogan: \emph{\(M\) gives
behavior; \(\Frame\) gives explanatory scope.}

The intervention model fixes the admissible counterfactual set
\(\Adm_M^{\Frame}(\tau)\) in one of three regimes:
\begin{description}[leftmargin=2.2cm,style=nextline]
  \item[Input-only] counterfactuals stay genuine machine behaviors,
        \(\Adm_M^{\Frame}(\tau)=\Lang(M)\cap\Lang(B_{\Frame})\). (Default; asks
        which environment behavior caused the effect.)
  \item[Output-intervention] outputs may be edited away from \(\lambda\):
        \(\Adm_M^{\Frame}(\tau)=\mathrm{Edit}_O^{\Frame}(\Lang(M))\cap
        \Lang(B_{\Frame})\), where \(\mathrm{Edit}_O^{\Frame}\) toggles outputs at
        a frame-permitted position set; a separate output/event-causality regime.
  \item[Strategy-structure] the system's commitments are re-opened:
        \(\Adm_M^{\Frame}(\tau)=\{\mathrm{Play}(\sigma',\rho)\mid\sigma'\in
        \Sigma_{\Frame},\ \rho\}\cap\Lang(B_{\Frame})\), with
        \[
          \Sigma_{\Frame}=\{\sigma'\mid \sigma'\text{ agrees with }\sigma
          \text{ on }V_0\setminus V_{\Frame}\}
        \]
        for a declared editable set \(V_{\Frame}\subseteq V_0\) of the synthesis
        game \(G_\varphi=(V,V_0,V_1,T,\Omega)\) and frozen strategy \(\sigma\)
        (we require \(\sigma\in\Sigma_{\Frame}\)). By default we take
        \(\sigma'\) to range over \emph{finite-memory} strategies on
        \(G_\varphi\) (sufficient for \(\omega\)-regular objectives on finite
        arenas); the witness search returns an accepting lasso \(\pi\) on a
        finite sub-arena, from which such a \(\sigma'\) exists by standard
        synthesis (Section~\ref{sec:rpg}). This regime requires inputs
        \((G_\varphi,\sigma,V_{\Frame})\) in addition to \(M\).
\end{description}
The \(\tau\)-argument is kept uniformly even though the first two regimes do not
depend on \(\tau\), because frame-relative interventions (the finite-horizon
TempoBench specialization of Section~\ref{sec:tempobench}) perturb \(\tau\) itself.

\subsection{Similarity to the actual trace}
\label{sec:similarity}

Counterfactual causality needs a notion of closeness, a preorder \(\preceq_\tau\)
where \(\tau'\preceq_\tau\tau''\) reads ``\(\tau'\) is at least as close to
\(\tau\) as \(\tau''\).'' Our canonical order is \emph{count-first}: let
\[
  \mathrm{Diff}_\tau(\tau')=\{t\mid\tau'(t)\neq\tau(t)\},\quad
  d(\tau,\tau')=|\mathrm{Diff}_\tau(\tau')|\in\Nat\cup\{\infty\},
\]
and define the primary key by \(d\) alone (fewer changed positions is closer;
\(\infty\) on top). Traces with an \emph{infinite} change set form a single
\(\preceq_\tau\)-maximal class we do not order internally; \(\Close\) (below) is
taken over \emph{finite-change} removals whenever any exist. Under this
\emph{count-only} order, \(\Close_{\tau,M}^{\Frame}(\neg C)\) is the full set of
minimum-count admissible \(\neg C\)-removals, and validity quantifies over
\emph{all} of them (effect must fail on each). This matches minimal-intervention
practice (Hamming distance on traces) and avoids an arbitrary tie-break changing
validity when two minimum-count removals differ only in \emph{which} positions
changed.

\begin{definition}[Optional profile refinement]
\label{def:profile-refine}
A frame may declare a total order \(\lhd_\Sigma\) on \(2^{I\cup O}\) and refine
\(\preceq_\tau\) lexicographically by difference profile \(\mathrm{prof}_\tau\)
(changed positions in order, each tagged with its new value). On finite-change
traces this yields a \emph{total} order with unique minimum \(\tau\) (centering)
and unique \(\Close\) (a single witness). Validity then depends on that tie-break;
reports must state \(\lhd_\Sigma\) or use count-only (default in
\(\Frame_{\mathrm{trace}}\) and \(\Frame_{\mathrm{TB}}\)).
\end{definition}

On finite-change traces the count-only order is antisymmetric among distinct
traces at fixed \(d\), well-founded (the count key is well-ordered), and has
unique minimum \(\tau\) among admissible traces. Lexicographic comparison of
\emph{infinite}-change profiles is \emph{not} well-founded, which is why
\(\Close\) is restricted to finite-change removals; if the only removals require
infinitely many edits, \(\Close\) is empty and the candidate fails non-vacuity.
This is a \textbf{scope boundary}, stated upfront: trace-level necessity (N1)
is confined to \emph{co-safety / finite-change} causes. An \(\omega\)-liveness
commitment whose removal requires infinitely many trace edits (e.g.\ removing
\(GF\,req\) from a lasso forces \(FG\,\neg req\)) has \(\Close=\emptyset\) at
the trace level and belongs to N2, where ``dropping a commitment'' is one
strategic deviation (Section~\ref{sec:rpg}). Neither N1 nor N2 alone is a
complete account of temporal explanation; the split reflects an inherent limit of
trace-level Lewis--Stalnaker necessity for liveness.

We reject two common alternatives on principled grounds: the first-difference
(Baire) order ties traces that differ only in \emph{when} the same minimal edit
occurs, and pure change-set \emph{inclusion} makes incomparable edits at the
same count --- both blur ``minimal intervention'' when several same-count
removals exist.

\begin{remark}[Closest removals for lasso counterfactuals]
\label{rem:lasso-close}
When \(\tau=uv^\omega\), \(M\) is finite-state, and the regime is input-only or
output-intervention, every admissible \(\neg C\)-trace is a bounded-period lasso
(a path in the finite product \(M\times A_{\neg C}\times A_E\)), and \(d\) is
finite exactly on perturbations agreeing with \(\tau\) on its tail. Under
count-only \(\preceq_\tau\), \(\Close_{\tau,M}^{\Frame}(\neg C)\) is the set of
minimum-count removals, found by BFS over that product; under the optional profile
refinement (Definition~\ref{def:profile-refine}) it is a unique shortest-edit
witness.
\end{remark}

\begin{remark}[Centering and the strategy-structure regime]
\label{rem:strat-centering}
Strategy-guarantee necessity (N2, Section~\ref{sec:two-notions}) is a
\emph{type-level} property of the controller \(\sigma\) --- ``can \(\sigma\) drop
a commitment and still guarantee \(E\)?'' --- and is centered on the actual
\emph{strategy} \(\sigma\), not on the actual trace \(\tau\). It therefore does
\emph{not} require trace-centering (Assumption~\ref{ass:wf}(ii)), and
Proposition~\ref{prop:rpg-sound} does not use it; the trace-level
Lewis--Stalnaker axiom is simply not the right notion at the strategy level.
Trace-centering matters only for the \emph{token-level} certificate
\(C^{\flat}\) (Proposition~\ref{prop:weakest}): if one wishes to run that
construction \emph{within} a strategy-structure frame (a separate, optional use),
the edited-vertex subset order ties all plays of \(\sigma\) with \(\tau\), so one
must refine it lexicographically (edited-vertex subset, then the trace-level
count-first order) to restore \(\tau\) as unique minimum. Absent that refinement,
Proposition~\ref{prop:weakest} is read only for the input-only/output regimes,
while the strategy-structure regime carries the validity content of
Proposition~\ref{prop:rpg-sound} (which needs no centering).
\end{remark}

\begin{assumption}[Well-formedness]
\label{ass:wf}
Fix \(M,\Frame,\tau\). (i)~\textbf{Actual-world admissibility:}
\(\tau\in\Adm_M^{\Frame}(\tau)\). (ii)~\textbf{Centering:} \(\preceq_\tau\) has
\(\tau\) as unique minimum (\(\tau\prec_\tau\tau'\) for admissible
\(\tau'\neq\tau\)). (iii)~\textbf{Well-foundedness (Limit Assumption):}
\(\preceq_\tau\) has no infinite strictly descending chain on
\(\Adm_M^{\Frame}(\tau)\), so every non-empty admissible set has minimal elements
and every element dominates a minimal one. Parts (i)--(ii) are Lewis--Stalnaker
centering; (iii) holds for the canonical order when \(\Close\) is restricted to
finite-change removals (Section~\ref{sec:similarity}), and is otherwise an explicit
hypothesis on \(\preceq_\tau\).
\end{assumption}

\subsection{Valid causes}

The closest admissible counterfactuals removing \(C\) are
\[
  \Close_{\tau,M}^{\Frame}(\neg C)=
  \min\nolimits_{\preceq_\tau}\{\tau'\in\Adm_M^{\Frame}(\tau)\mid\tau'\models\neg C\}
\]
(over finite-change removals; Section~\ref{sec:similarity}; under count-only
\(\preceq_\tau\) this is the full minimum-count set). The valid-cause set is
\[
\begin{aligned}
  \mathsf{Causes}_{M,\Frame,\tau}(E_{\Frame})=\{C\in\mathcal{K}_{\Frame}\mid{}&
  \tau\in\Lang(C),\ \tau\in\Lang(E_{\Frame}),\\
  &\Close_{\tau,M}^{\Frame}(\neg C)\neq\emptyset,\\
  &\Close_{\tau,M}^{\Frame}(\neg C)\cap\Lang(E_{\Frame})=\emptyset\}.
\end{aligned}
\]
The third (non-vacuity) clause is essential: without it any \(C\) that no
admissible trace violates --- a tautology, or any invariant of \(M\) under an
input-only frame --- would pass vacuously. The fourth clause is counterfactual
dependence on \emph{every} closest removal (under count-only) or the unique
closest removal (under profile refinement).

\paragraph{Strategy-level causes need a re-opening frame.}
\begin{proposition}
\label{prop:strategy-frame}
Let \(C_M\) be an invariant of \(M\) (\(\Lang(M)\subseteq\Lang(C_M)\)). Under an
input-only frame, \(C_M\notin\mathsf{Causes}_{M,\Frame,\tau}(E_{\Frame})\) for
every \(\tau,E_{\Frame}\).
\end{proposition}
\begin{proof}
\(\Adm_M^{\Frame}(\tau)\subseteq\Lang(M)\subseteq\Lang(C_M)\), so no admissible
trace violates \(C_M\); \(\Close_{\tau,M}^{\Frame}(\neg C_M)=\emptyset\) and
non-vacuity fails.
\end{proof}
Thus an invariant is \emph{vacuous under input-only} counterfactuals; testing it
requires re-opening the machine's commitments, the principled regime for which
is strategy-structure intervention (the reverse parity game).

\subsection{Anti-circularity}

\(\mathcal{K}_{\Frame}\) must be constrained, or the test admits circular
``causes'' that re-assert the effect. The key distinction is between
\emph{causing \(E\) through \(M\)} (what a genuine cause does, and must be allowed)
and \emph{restating \(E\)} (asserting the effect at its own evaluation positions,
which is circular). The condition bans the latter without touching the former.

\begin{assumption}[Anti-circularity]
\label{ass:anticirc}
Every \(C\in\mathcal{K}_{\Frame}\) is \emph{effect-disjoint}:
\begin{enumerate}[label=(\roman*)]
  \item \textbf{No restatement (semantic, point/bounded effects).} For a point
        effect \(E_{\Frame}=X^k o\), \(C\) does not constrain the value of \(o\) at
        position \(k\): \(\Lang(C)\) is invariant under flipping \(o\) at position
        \(k\) (i.e.\ \(\tau'\in\Lang(C)\iff\tau'[k{:}o{\mapsto}\neg o]\in\Lang(C)\),
        where the bracket flips only \(o\) at time \(k\)). Equivalently, in the
        masked vocabulary that hides \(o@k\), \(C\) is well-defined. This is the
        clean, semantic form of ``\(C\) does not assert the effect at its own
        position''; it admits a cooldown cause ``\(o\) at \(t=5\)'' for effect
        ``\(o\) at \(t=10\)'' (\(t\neq k\)) and any inputs, while barring
        \(X^k o\), \(X^k\neg o\), and trivial re-encodings such as
        \(X^k(o\lor\bot)\). For a bounded conjunctive effect the masking is applied
        at each constrained position. \emph{For a general \(\omega\)-regular
        effect} (\(GF\,grant\), \(G(req\to F\,grant)\), parity acceptance) there is
        no single ``effect position'' to mask; guard (ii) and a declared
        \emph{effect cone} (Assumption~\ref{ass:cone}, mandatory) apply instead.
  \item \textbf{No pure-logic entailment (all effects).} \(\Lang(C)\not\subseteq
        \Lang(E_{\Frame})\), decidable by emptiness of \(\Lang(C)\cap
        \overline{\Lang(E_{\Frame})}\) (an \(\omega\)-regular inclusion test). The
        test is \emph{not} relativized to \(\Adm_M^{\Frame}(\tau)\): a candidate
        that forces \(E_{\Frame}\) only \emph{through \(M\)} is \emph{not} excluded
        (this is what reactive causes do); only one that forces \(E_{\Frame}\) by
        pure logic --- already a logical strengthening of the effect --- is. This
        bars logical entailment and point-position restatement; it does \emph{not}
        exclude system-relative aliases or near-restatements that pass (ii) alone
        (hence the mandatory effect cone below).
\end{enumerate}
We call this ``precise and decidable'' in the sense that (i) is a decidable
semantic invariance check for point/bounded effects, (ii) is a decidable
\(\omega\)-regular inclusion test for all effects, and the effect cone is a
declared finite AP subset; we do \emph{not} claim a uniform syntactic
``effect-position'' notion for arbitrary temporal effects.
\end{assumption}

\begin{assumption}[Effect cone (mandatory for general effects)]
\label{ass:cone}
For every non-point effect \(E_{\Frame}\), the frame declares a finite
\(\AP\)-subset \(\mathsf{Cone}(E_{\Frame})\) of propositions that may appear in
effect-shaped subformulas (outputs and acceptance atoms tied to \(E\)). Every
\(C\in\mathcal{K}_{\Frame}\) must satisfy \(\AP(C)\cap\mathsf{Cone}(E_{\Frame})
=\emptyset\) unless \(C\) is explicitly marked, in \(\Frame_{\mathrm{strat}}\), as
an internal-state observable (e.g.\ \(\mathit{at\_}s\)) or a declared
system-guarantee invariant. This blocks vacuous near-restatements that name the
effect atom --- e.g.\ a bare recurrence \(GF\,grant\), or a guarded re-assertion of
\(grant\) --- for effect \(GF\,grant\) when \(\mathsf{Cone}=\{grant\}\), while still
allowing input causes and, under a declared regime, strategy-level state
observables and system-guarantee invariants (e.g.\ the controller commitment
\(G(req\to F\,grant)\), admitted only by such an explicit declaration, never at the
trace level).
\end{assumption}

This is what makes the worked causes admissible. For the general effect
\(F\,grant\), guard (i) is silent and (ii) is operative: the input cause \(req\)
satisfies \(\Lang(req)\not\subseteq\Lang(F\,grant)\) (a \(req\)-at-0, grant-free
trace witnesses this), so it is admitted --- yet \(req\) does force \(F\,grant\) on
\(\Lang(M)\), which is precisely why it is a cause. The TempoBench cause
\(\Gamma^\star\) (Section~\ref{sec:tempobench}), pinning \(o\) at \(k\) through the
functional machine, is admitted similarly (for the point effect \(X^k o\), guard
(i) masks only \(o@k\), which \(\Gamma^\star\), being over inputs, does not touch).
The invariant \(G(req\rightarrow F\,grant)\) names the cone atom \(grant\), so when
\(\mathsf{Cone}(F\,grant)=\{grant\}\) it is admitted \emph{only} as a declared
system-guarantee invariant under \(\Frame_{\mathrm{strat}}\) (never at the trace
level); given that declaration it is not a logical restatement either, since a
request-free trace satisfies it vacuously while failing the effect. \emph{Excluded} are logical
restatements: \(X\,grant\) and \(req\land X\,grant\) entail \(F\,grant\) (fail
(ii)); \(X^k o\) and \(X^k(o\lor\bot)\) for point effect \(X^k o\) fail (i).

\begin{remark}[Strategy-level state observables]
\label{rem:strat-obs}
Under \(\Frame_{\mathrm{strat}}\), candidates may use \(\mathit{at\_}s\) atoms
(exposing internal state). Anti-circularity (ii) still requires
\(\Lang(C)\not\subseteq\Lang(E_{\Frame})\); coincidence of \(C\) with \(E\) on
\(\Lang(M)\) is machine entailment, not logical restatement. Such observables are
admissible when the explanatory question is about a controller \emph{commitment},
not a trace event (Section~\ref{sec:liveness-example}).
\end{remark}

\subsection{Minimal cause and the preference order}

\begin{definition}[Minimal temporal cause]
\label{def:mincause}
\(C^\star\in\mathcal{K}_{\Frame}\) is a minimal temporal cause iff
\(C^\star\in\mathsf{Causes}_{M,\Frame,\tau}(E_{\Frame})\) (validity) and no
\(C'\in\mathsf{Causes}_{M,\Frame,\tau}(E_{\Frame})\) has \(C'\triangleleft C^\star\)
(minimality). When \(\triangleleft\) is partial, several incomparable minimal
causes may exist --- ``\emph{a} minimal cause,'' not ``the.'' We assume
\(\triangleleft\) well-founded (automatic when \(\mathcal{K}_{\Frame}\) is finite).
\end{definition}

Validity already encodes necessity; \(\triangleleft\) selects \emph{among}
necessary causes. Taking permissiveness (larger \(\Lang(C)\)) as the
\emph{primary} key is defective: if \(A,B\) are each valid then \(A\lor B\) is
weaker, so a permissiveness-first order drifts to ever-larger disjunctions and
would even prefer \(F\,req\) to \(req\). We therefore make \emph{syntactic
readability primary and permissiveness only a tie-break}; the canonical order is
\[
  (\text{smaller formula size},\ \text{smaller temporal depth},\
   \text{more permissive (larger }\Lang(C)\text{)}),
\]
lexicographic. For ``size'' and ``depth'' to be reproducible we fix a candidate
grammar (the readable LTL grammar of Section~\ref{sec:construction}) and a normal
form (negations on literals, fixed \(\land,\lor\) ordering, no redundant
\(X\)-nesting), measuring size as syntax-tree node count; the grammar and measure
are part of the frame. We do not claim this priority is uniquely correct --- it is
a declared design choice; the frame may substitute another well-founded order. The
default has three consequences: (i)~it avoids overfitting (large conjunctions lose
on size); (ii)~it avoids disjunctive drift (\(A\lor B\) is larger, chosen only
when no smaller cause exists --- exactly overdetermination,
Section~\ref{sec:overdetermination}); (iii)~the degenerate extensionally maximal
object below has no small readable formula. We require of \(\mathcal{K}_{\Frame}\)
that the trace-exclusion artifact \(A\setminus\{f_0\}\) be inexpressible (it holds
for the readable grammar/template families). If the necessary condition
lies outside \(\mathcal{K}_{\Frame}\) (e.g.\ \(a\lor b\) when only literals are
admitted), the method reports \(\mathsf{Causes}=\emptyset\); the remedy is to
enlarge \(\mathcal{K}_{\Frame}\), not to return a spurious singleton.

\begin{theorem}[Decidability for fixed finite candidate space]
\label{thm:finite-K}
Fix \((M,\Frame,\tau,E_{\Frame})\) with \(|\mathcal{K}_{\Frame}|<\infty\). For
each \(C\in\mathcal{K}_{\Frame}\), membership in
\(\mathsf{Causes}_{M,\Frame,\tau}(E_{\Frame})\) is decidable (N1: non-emptiness
plus shortest-edit search, Remark~\ref{rem:lasso-close}; anti-circularity is
decidable by Assumption~\ref{ass:anticirc}). If \(\triangleleft\) is
well-founded on \(\mathcal{K}_{\Frame}\), a \(\triangleleft\)-minimal valid
\(C^\star\) exists and is found by finite search. This is the precise ``compute a
minimal cause'' claim; unbounded LTL grammar search is only decidable under an
explicit size bound (Section~\ref{sec:complexity}).
\end{theorem}

\subsection{Two distinct necessity notions}
\label{sec:two-notions}

We state \emph{once} the central distinction (it recurs throughout and we
cross-reference here rather than repeat it).

\begin{center}
\small
\begin{tabular}{@{}l|l|l@{}}
\hline
 & \textbf{N1 (trace-existential)} & \textbf{N2 (strategy-guarantee)} \\
\hline
Inputs & \(M,\Frame,E,\tau\) & \(G_\varphi,\sigma,V_{\Frame},\Frame,E\) \\
Centering & actual trace \(\tau\) & actual strategy \(\sigma\) \\
Question & closest \(\neg C\) trace fails \(E\)? &
  can \(\sigma\) drop \(C\) and still guarantee \(E\)? \\
Machinery & one-player shortest-edit + non-emptiness &
  reverse parity game + re-imposed objective \\
Liveness & fails if \(\Close=\emptyset\) & primary regime \\
\hline
\end{tabular}
\end{center}

\begin{description}[leftmargin=2.2cm,style=nextline]
  \item[Trace-existential necessity (N1)] the literal \(\mathsf{Causes}\)
        definition above: every closest admissible \(\neg C\)-trace (under
        count-only \(\preceq_\tau\), or the unique one under profile refinement)
        must fail \(E\). Deciding N1 for fixed \(C\) is non-emptiness of
        \(\Adm\cap\Lang(\neg C)\) over finite-change removals plus a
        shortest-edit search in the product \(M\times A_{\neg C}\times A_E\)
        (Remark~\ref{rem:lasso-close}), then evaluating \(E\) on the result --- not
        emptiness alone. Token-level: ``this execution.''
  \item[Strategy-guarantee necessity (N2)] the two-player, strategy-level
        question: can the system \emph{drop the commitment} \(C\) and still
        \emph{guarantee} \(E\) against every \emph{admissible} environment?
        ``Admissible'' means satisfying the background \(B_{\Frame}\) (e.g.\ a
        fairness assumption); quantifying over \emph{all} environments, including
        \(B_{\Frame}\)-violating ones, would let an unfair environment trivially
        falsify any liveness commitment, so both clauses are relativized to
        \(B_{\Frame}\). Formally, the frozen \(\sigma\) is assumed to secure the
        guarantee and to satisfy \(C\) (the actuality requirement at the strategy
        level); \(C\) is \emph{strategy-necessary} iff \textbf{no}
        \(\sigma'\in\Sigma_{\Frame}\) both
        \begin{itemize}
          \item[(i)] \emph{secures} \(E\) on admissible plays:
                \(\forall\rho.\ \mathrm{Play}(\sigma',\rho)\models
                (B_{\Frame}\rightarrow E)\), and
          \item[(ii)] is \emph{not} \(C\)-committed on an admissible play:
                \(\exists\rho.\ \mathrm{Play}(\sigma',\rho)\models
                B_{\Frame}\wedge\neg C\).
        \end{itemize}
        Type-level: ``this controller's guarantee.'' (When \(B_{\Frame}=\True\),
        (i) is \(\forall\rho.\,E\) and (ii) is \(\exists\rho.\,\neg C\), as in the
        un-relativized form.) N2 is centered on the actual \emph{strategy}
        \(\sigma\), not on \(\tau\): it is a property of the controller, so the
        trace-level centering of \(\tau\) plays no role here (cf.\
        Remark~\ref{rem:strat-centering}).
\end{description}
The reverse parity game (Section~\ref{sec:rpg}) decides N2; N1 is the one-player
check. We do \emph{not} claim a single parity game decides N1.

\subsection{The certificate \texorpdfstring{\(C^{\flat}\)}{C-flat}}
\label{sec:certificate}

Write \(A=\Adm_M^{\Frame}(\tau)\) and \(F=\{\tau'\in A\mid\tau'\not\models
E_{\Frame}\}\) (admissible effect-failures), and
\(\uparrow_{\preceq_\tau}(F)=\{\tau''\mid\exists\tau'\in F.\ \tau'\preceq_\tau
\tau''\}\) (here ``above'' means \emph{farther} from \(\tau\)). Let \(C^{\flat}\)
be \emph{any} \(\omega\)-regular property with
\(\Lang(C^{\flat})\cap A=A\setminus\uparrow_{\preceq_\tau}(F)\); this pins
\(C^{\flat}\) only on \(A\), so ``\(C^{\flat}\)'' names an equivalence class, not a
unique formula.

\begin{proposition}[Nearest-failure-anchored certificate]
\label{prop:weakest}
Suppose Assumption~\ref{ass:wf} holds, \(\preceq_\tau\) is well-founded on \(A\),
\(\tau\models E_{\Frame}\), and \(F\neq\emptyset\). Then:
\begin{enumerate}[label=(\alph*)]
  \item \(C^{\flat}\) satisfies the four valid-cause \emph{clauses}; hence it is a
        valid cause \emph{provided} \(C^{\flat}\in\mathcal{K}_{\Frame}\). As a
        distance-defined object \(C^{\flat}\) need not lie in \(\mathcal{K}_{\Frame}\)
        (it may be unreadable or fail effect-disjointness), so we do not assert
        \(C^{\flat}\in\mathsf{Causes}\) unconditionally; it is a witness, not the
        returned answer \(C^\star\);
  \item \(\Close_{\tau,M}^{\Frame}(\neg C^{\flat})=\min_{\preceq_\tau}(F)\);
  \item among valid causes whose closest removals are the globally nearest
        effect-failures, \(C^{\flat}\) is the \emph{smallest-language} one:
        \(\Lang(C^{\flat})\cap A\subseteq\Lang(C)\cap A\) for every such \(C\).
        Equivalently --- and this is not a contradiction --- its admissible models
        \(\Lang(C^{\flat})\cap A=A\setminus\uparrow(F)\) form the \emph{largest}
        failure-free neighborhood of \(\tau\), because its complement
        \(\uparrow(F)\) is the \emph{smallest} failure-containing up-set; the cause
        is language-minimal precisely because the failures it must admit are pushed
        as far out as possible.
\end{enumerate}
\end{proposition}
\begin{proof}
Within \(A\), \(\neg C^{\flat}\equiv\uparrow(F)\). For (b): by antisymmetry and
well-foundedness \(\min(\uparrow(F))=\min(F)\) (if \(m\in\min(\uparrow F)\) then
some \(f\in F\) has \(f\preceq_\tau m\), minimality gives \(m\preceq_\tau f\), and
antisymmetry \(m=f\in F\); the converse is similar). For (a): \(\min(F)\neq
\emptyset\) by the Limit Assumption (non-vacuity); every closest removal lies in
\(F\), failing \(E_{\Frame}\); \(\tau\models E_{\Frame}\) by hypothesis and
\(\tau\models C^{\flat}\) since centering gives \(\tau\notin\uparrow(F)\). For (c):
let \(C\) be valid with \(\Close_{\tau,M}^{\Frame}(\neg C)=\min_{\preceq_\tau}(F)\).
If \(\tau'\in\Lang(C^{\flat})\cap A\) but \(\tau'\notin\Lang(C)\), then
\(\tau'\models\neg C\), so some \(\tau''\in\Close_{\tau,M}^{\Frame}(\neg C)\)
satisfies \(\tau''\preceq_\tau\tau'\) (well-founded minimality). Since
\(\Close_{\tau,M}^{\Frame}(\neg C)=\min(F)\subseteq F\), we have \(\tau''\in F\),
hence \(\tau'\in\uparrow(F)\), contradicting
\(\tau'\in A\setminus\uparrow(F)=\Lang(C^{\flat})\cap A\).
\end{proof}

\begin{remark}[The extensionally maximal cause is the wrong target]
\label{rem:weakest-degenerate}
The \emph{extensionally maximal} (largest-language) valid cause --- we avoid
``weakest'' to prevent clash with Dijkstra's weakest precondition --- is
degenerate: removing a single nearest failure \(f_0\) gives the valid cause
\(A\setminus\{f_0\}\), maximal yet non-explanatory and arbitrary in \(f_0\). The
\emph{answer} is the \(\triangleleft\)-minimal readable \(C^\star\); \(C^{\flat}\)
is a \emph{theoretical certificate} (smallest anchored language on \(A\), generally
unreadable, not necessarily in \(\mathcal{K}_{\Frame}\)); it lower-bounds
\(\Lang(C)\cap A\) among anchored valid causes but does not by itself yield
\(C^\star\) --- search still enumerates \(\mathcal{K}_{\Frame}\) or runs CEGIS
(Section~\ref{sec:construction}). The extensionally maximal artifact is excluded
(unreadable and large). These three objects --- \(C^\star\) (answer), \(C^{\flat}\)
(certificate anchor), extensionally maximal (excluded) --- must not be conflated. \(C^{\flat}\) stays
\(\omega\)-regular when \(\uparrow_{\preceq_\tau}(F)\) is (true for the count-first
order via Remark~\ref{rem:lasso-close} on lassos; otherwise an explicit
automaticity hypothesis). We do not claim the order \emph{relation} is automatic in
general.
\end{remark}

A useful output is the triple \((C^\star,W,\Pi)\): the cause \(C^\star\); the
preserved region \(W\) (accepting region of \(P=M\times A_E\) in one-player regimes,
or \(W_\exists\) in strategy-structure); and counterfactual witnesses \(\Pi\)
(one per rejected competitor \(C'\), each a lasso
\(\pi_{C'}\in\Close_{\tau,M}^{\Frame}(\neg C')\cap\Lang(E_{\Frame})\) from the
emptiness/shortest-edit check or Proposition~\ref{prop:rpg-sound}).

\subsection{Canonical frames}
\label{sec:canonical-frames}

To answer the objection that a free choice of \(\Frame\) injects the cause, we fix
canonical defaults; results are reported against one unless a deviation is
justified.
\begin{description}[leftmargin=2.4cm,style=nextline]
  \item[\(\Frame_{\mathrm{trace}}\)] \(B_{\Frame}=\True\) (or stated environment
        assumptions), input-only, count-only \(\preceq_\tau\) (Definition~\ref{def:profile-refine}:
        no profile tie-break unless declared), effect-disjoint candidates over
        \(I\cup O\) with declared \(\mathsf{Cone}(E_{\Frame})\), and
        \(\triangleleft=\) (size, depth, permissiveness). Default for one execution;
        the toy example uses optional profile refinement for a unique witness.
  \item[\(\Frame_{\mathrm{strat}}\)] as above but strategy-structure over stated
        \((G_\varphi,\sigma,V_{\Frame})\) with edited-vertex order refined by trace
        distance (Remark~\ref{rem:strat-centering}), internal-state observables
        allowed (Remark~\ref{rem:strat-obs}), tested via Section~\ref{sec:rpg}.
  \item[\(\Frame_{\mathrm{TB}}\)] the finite-horizon, point-effect, input-only
        frame of Section~\ref{sec:tempobench}.
\end{description}
Two safeguards: every component but candidate vocabulary, grammar, similarity
tie-break, and effect cone is fixed by the regime; and under a \emph{fully
specified} canonical frame (including those declarations)
\(\mathsf{Causes}_{M,\Frame,\tau}(E_{\Frame})\) is a function of
\((M,\tau,E_{\Frame})\) alone. Frame-relativity means results are conditional on
that declaration, as background conditions are in ordinary causal explanation.
Determinacy requires stating vocabulary, grammar, count-only vs.\ profile
\(\preceq_\tau\), and \(\mathsf{Cone}(E_{\Frame})\) for non-point effects.

\section{Trace-Level Construction and Toy Example}
\label{sec:construction}

\subsection{The trace-level procedure (one-player)}

Given \((M,\Frame,E,\tau)\), the trace-level test (N1) is: (1)~compile \(E\) into
a deterministic parity automaton \(A_E=(Q,q_0,2^{I\cup O},\delta_E,\Omega)\) and
form the frozen product \(P=M\times A_E\) (a \emph{one-player} arena: \(M\) is
fixed, so only the environment moves); (2)~for a candidate \(C\), build
\(A_{\neg C}\) and check
\[
  \Close_{\tau,M}^{\Frame}(\neg C)\neq\emptyset
  \quad\text{and}\quad
  \Close_{\tau,M}^{\Frame}(\neg C)\cap\Lang(E)=\emptyset ,
\]
an emptiness check plus shortest-edit search (Remark~\ref{rem:lasso-close}). To
\emph{search} for a \(\triangleleft\)-minimal valid cause when
\(|\mathcal{K}_{\Frame}|<\infty\), enumerate candidates (Theorem~\ref{thm:finite-K});
otherwise use bounded grammar/template search with CEGIS:
\[
  C ::= p \mid \neg C \mid C\land C \mid X C \mid F C \mid G C \mid C\,U\,C,
\]
or templates \(G(a\to F b),F(a\land X b),G(a\to X b),GF a,FG a\), proposing \(C\),
checking N1, and refining from witnesses \(\Pi\). The certificate \(C^{\flat}\)
(Section~\ref{sec:certificate}) characterizes the anchored language lower bound but
does not replace this search. Strategy-level minimality uses dominion subset search
(Section~\ref{sec:rpg}). The returned formula need not be a subformula of
\(\varphi\); we translate readable LTL from \(\mathcal{K}_{\Frame}\) when needed.

\subsection{Toy example}
\label{sec:toy}

Let \(\varphi=G(req\rightarrow F\,grant)\) and let \(M\) have states
\(idle,pending\) with initial state \(s_0=idle\), specified completely:
\[
\begin{array}{rclcl}
  \delta(idle,req)&=&pending, &\quad& \lambda(idle,\cdot)=\emptyset,\\
  \delta(idle,\neg req)&=&idle, &\quad& \lambda(pending,\cdot)=\{grant\},\\
  \delta(pending,\cdot)&=&idle. & &
\end{array}
\]
A request in \(idle\) moves to \(pending\) next; in \(pending\) the machine emits
\(grant\) and returns to \(idle\); a request in \(pending\) is \emph{not} latched.
The machine is effectively Moore (\(grant\) holds exactly in \(pending\)). The
actual trace is \(\{req\},\{grant\},\emptyset,\emptyset,\ldots\) and the effect is
\(E=F\,grant\).

\paragraph{The minimal trace-level cause is \(req\).} Use \(\Frame_{\mathrm{trace}}\)
with count-only \(\preceq_\tau\) (\(\Adm=\Lang(M)\)). Every minimum-count removal
sets \(\neg req\) at time \(0\): the machine stays in \(idle\), never enters
\(pending\), never grants; \(d=2\) and \(F\,grant\) fails. No \(\neg req\)-removal
of count \(\le2\) preserves the effect (emitting \(grant\) needs a \(req\) in
\(idle\) at some \(j\ge1\); the cheapest such trace has \(d=3\)). All minimum-count
removals fail \(F\,grant\); so \(req\) is non-vacuous and necessary --- a valid
cause. (Under profile refinement the unique closest removal is the \(d=2\) trace
above; under Baire or inclusion orders \(req\) would spuriously fail.) The candidate
\(req\land X\,grant\) is barred by Assumption~\ref{ass:anticirc}(ii) and loses on
size. Hence \(C_\tau=req\).

\paragraph{Why \(req\) can feel uninformative.} That ``a request alone does not
explain \emph{why} the machine granted'' is a critique of the input-only
\emph{frame}, which fixes the machine's response rule as background: relative to
that, the only thing that had to hold for the grant was that a request arrived. To
make the response rule a candidate, move to a strategy frame. The strategy-level
target is \(C_M=G(req\rightarrow F\,grant)\), or, exposing the state,
\(G((\mathit{at\_idle}\land req)\rightarrow X\,pending)\land G(pending\rightarrow
grant)\) (the first conjunct must be guarded by \(\mathit{at\_idle}\), since a
request in \(pending\) is not latched, so unguarded \(G(req\to X\,pending)\) is not
an invariant). These entail \(G((\mathit{at\_idle}\land req)\rightarrow X\,grant)\),
hence \(G(req\rightarrow F\,grant)\) on traces whose requests arrive only in
\(idle\). By Proposition~\ref{prop:strategy-frame}, \(C_M\) is testable only under
a strategy frame; it names the cone atom \(grant\), so when
\(\mathsf{Cone}(F\,grant)=\{grant\}\) it is admissible only because
\(\Frame_{\mathrm{strat}}\) declares it a system-guarantee invariant
(Assumption~\ref{ass:cone}), not at the trace level. The two answers --- \(req\) (trace) and
\(G(req\to F\,grant)\) (strategy) --- are outputs of two different frames, not
competitors.

\section{The Reverse Parity Game (Strategy Level)}
\label{sec:rpg}

Many temporal effects (\(GF\,a\), \(FG\,stable\), \(G(req\to F\,ack)\)) cannot be
explained by a single finite event; the parity condition is needed to reason about
recurrence. But it is needed \emph{only} for N2: in the input-only and
output-intervention regimes the trace test is one-player
(Section~\ref{sec:two-notions}). This section gives the genuine two-player object
for N2. (``Reverse'' here names the \emph{explanatory} direction --- we run the
synthesis machinery backwards, from a frozen controller to the conditions it
needed --- not a retrograde or backward-time analysis; the game itself is a
standard forward parity game on a partial-commitment arena.) Fix the synthesis game \(G_\varphi=(V,V_0,V_1,T,\Omega)\), the frozen
\(\sigma\) implemented by \(M\), the candidate \(C\) with deterministic parity
monitor \(A_{\neg C}\), the effect monitor \(A_E\), and the editable vertex set
\(V_{\Frame}\subseteq V_0\) defining \(\Sigma_{\Frame}\).

\begin{definition}[Reverse parity game]
\label{def:rpg}
\(G^{\mathrm{rev}}_{M,\Frame,C,E}=(V^r,V^r_0,V^r_1,T^r,\Omega^r)\) has vertices
\(V^r=V\times Q_B\times Q_E\times Q_{\neg C}\), pairing a game vertex with the
states of deterministic monitors \(A_B\) (for the background \(B_{\Frame}\)),
\(A_E\), and \(A_{\neg C}\). It is \emph{partitioned} so that \(V^r_0\)
(\(\exists\), the re-opened system) are vertices over \(V_{\Frame}\), where
\(\exists\) may take any \(G_\varphi\)-move, and \(V^r_1\) (\(\forall\)) are all
others: vertices over \(V_1\) (environment) and over the frozen
\(V_0\setminus V_{\Frame}\) (single-successor, assigned to \(\forall\) to keep a
genuine partition while leaving play unaffected). Edges carry the synthesis-game
label \(\ell:T\to 2^{I\cup O}\); \(T^r\) updates \(V\) by \(T\) and each monitor by
its transition function on \(\ell\). The \(\exists\)-objective is the
assume-guarantee condition \(B_{\Frame}\rightarrow E\): \(\exists\) wins \(\pi\)
iff \(\pi|_{I\cup O}\models B_{\Frame}\rightarrow E\). This is \(\omega\)-regular,
so \(\Omega^r\) is a parity index for it and \(G^{\mathrm{rev}}\) is a genuine
parity game; \(A_B\) and \(A_{\neg C}\) are also read off as \emph{auxiliary}
acceptance conditions used in the witness search below.
\end{definition}

N2 is a \emph{conjunctive} strategy-existence question --- \(\exists\sigma'\) that
secures \(B_{\Frame}\rightarrow E\) (\(\forall\rho\)) \emph{and} admits an
admissible \(\neg C\) play (\(\exists\rho.\ B_{\Frame}\wedge\neg C\)) --- and
ordinary parity solving returns one winning strategy and the winning region, not
the set of all winning strategies. \textbf{Caution:} for a parity/B\"uchi
objective, staying inside the winning region forever does \emph{not} imply winning
(a play may remain in \(W_\exists\) yet violate the parity condition), so we do
\emph{not} use a ``trap'' multi-strategy of ``all moves staying in \(W_\exists\)'';
we re-impose the objective explicitly in the witness search.

\begin{proposition}[Reverse-game soundness for N2]
\label{prop:rpg-sound}
Assume the frozen \(\sigma\in\Sigma_{\Frame}\) secures \(B_{\Frame}\rightarrow E\)
and satisfies \(C\) (the strategy-level actuality requirement). Let \(W_\exists\)
be the \(\exists\)-winning region of \(G^{\mathrm{rev}}_{M,\Frame,C,E}\) for
objective \(B_{\Frame}\rightarrow E\), and let \(\mathcal{A}[W_\exists]\) be the
sub-arena obtained by restricting \(G^{\mathrm{rev}}\) to \(W_\exists\) (keeping
only edges that stay in \(W_\exists\)). Then \(C\) is \emph{not}
strategy-necessary (N2 fails) iff \(\mathcal{A}[W_\exists]\) contains a play
\(\pi\) with
\[
  \pi|_{I\cup O}\models B_{\Frame}\wedge E\wedge\neg C .
\]
Equivalently, \(C\) is strategy-necessary iff the product
\(\mathcal{A}[W_\exists]\times A_B\times A_E\times A_{\neg C}\) has no accepting
lasso for the conjunction \(B_{\Frame}\wedge E\wedge\neg C\) --- a one-player
\(\omega\)-regular non-emptiness check (not mere reachability of a vertex).
\end{proposition}
\begin{proof}
Let \(\Pi_{\mathrm{win}}\) be the plays consistent with \emph{some}
\(\sigma'\in\Sigma_{\Frame}\) that secures \(B_{\Frame}\rightarrow E\). We claim
\(\Pi_{\mathrm{win}}=\{\pi:\pi\text{ stays in }W_\exists\text{ and }\pi\models
B_{\Frame}\rightarrow E\}\). (\(\subseteq\)) A securing strategy keeps every play
in \(W_\exists\) and wins it, so its plays satisfy \(B_{\Frame}\rightarrow E\).
(\(\supseteq\)) Given such a \(\pi\), we may take \(\pi=uv^\omega\) to be an
accepting \emph{lasso}: \(\mathcal{A}[W_\exists]\) is finite and the winning
condition \(\omega\)-regular, so whenever a qualifying play exists a lasso one
does (this is the lasso/non-emptiness fact, standard for deterministic
\(\omega\)-automata over a finite arena). Define \(\sigma'\) to follow the
(ultimately periodic) moves of \(\pi\) while the history matches \(\pi\), and a
fixed \emph{finite-memory} winning strategy from \(W_\exists\) after any
environment deviation: every \(\exists\)-move on \(\pi\) goes to \(W_\exists\)
(\(\pi\) stays there), and the environment cannot escape \(W_\exists\). Following a
lasso needs only finite memory (the period), and finite-memory winning strategies
exist for parity objectives on finite arenas, so \(\sigma'\) is finite-memory;
hence \(\sigma'\in\Sigma_{\Frame}\) (which is the finite-memory class by default,
Section~\ref{sec:formal}), \(\sigma'\) secures \(B_{\Frame}\rightarrow E\), and
\(\pi\) is one of its plays. Now \(C\) is not strategy-necessary iff some
\(\sigma'\) secures \(B_{\Frame}\rightarrow E\) (i) and admits a play with
\(B_{\Frame}\wedge\neg C\) (ii), i.e.\ iff \(\Pi_{\mathrm{win}}\) contains a play
with \(B_{\Frame}\wedge\neg C\). By the claim such a play stays in \(W_\exists\)
and satisfies \(B_{\Frame}\rightarrow E\); with \(B_{\Frame}\) this gives \(E\), so
it satisfies \(B_{\Frame}\wedge E\wedge\neg C\) and lies in
\(\mathcal{A}[W_\exists]\). Conversely any \(\mathcal{A}[W_\exists]\)-play with
\(B_{\Frame}\wedge E\wedge\neg C\) satisfies \(B_{\Frame}\rightarrow E\) and stays
in \(W_\exists\), hence lies in \(\Pi_{\mathrm{win}}\) and witnesses (i)+(ii). The
lasso/non-emptiness characterization is standard for deterministic
\(\omega\)-automata over a finite arena.
\end{proof}

This decision is sound \emph{and} complete and uses only (a)~a two-player parity
solve for the winning region \(W_\exists\) and (b)~a one-player
\(\omega\)-regular non-emptiness check on \(\mathcal{A}[W_\exists]\). It does
\emph{not} rely on a maximal-permissive-strategy construction, sidestepping the
fact that for parity objectives the ``moves staying in \(W_\exists\)''
multi-strategy is \emph{not} winning (a Bernet--Janin--Walukiewicz most-permissive
\emph{winning} strategy needs memory); the crucial point is that the objective
\(B_{\Frame}\rightarrow E\) is re-imposed in the witness search, so non-winning
plays that merely stay in \(W_\exists\) are excluded. The input-only/output
regimes are the degenerate case \(V_{\Frame}=\emptyset\) (\(V^r_0=\emptyset\)):
the game is one-player and the question collapses to the trace check of
Section~\ref{sec:construction} (N1). Minimizing \(V_{\Frame}\) is a separate
optimization (Section~\ref{sec:complexity}).

\paragraph{From commitments to a trace cause.} If \(C\) is strategy-necessary
(no \(\mathcal{A}[W_\exists]\)-play satisfies \(B_{\Frame}\wedge E\wedge\neg C\)),
the witness is the set \(D_C\subseteq V_{\Frame}\) of editable vertices whose
\(\sigma\)-move, if re-opened, would let \(\exists\) secure
\(B_{\Frame}\rightarrow E\) via a \(\neg C\)-deviation, with
transitions \(T_C=\{(v,\sigma(v))\mid v\in D_C\}\). For a \emph{safety} commitment
the induced cause is \(C\equiv G\bigwedge_{v\in D_C}(\mathit{at\_}v\to
\mathit{move}_\sigma(v))\) (negation \(F\bigvee(\ldots)\)); for a \emph{liveness}
commitment, e.g.\ \(C=GF\,\mathit{at\_serve}\) below, it is
\(GF\bigvee_{v\in D_C}\mathit{at\_}v\). Breaking \(T_C\) is exactly what a
\(\neg C\)-deviation does.

\paragraph{Minimal strategy cause.} The strategy level so far has only
\emph{validity} (N2). A \emph{minimal strategy cause} is a \(\subseteq\)-minimal
editable set \(V_{\Frame}\) (equivalently a minimal dominion \(D_C\)) under which
\(C\) is strategy-necessary. Unlike the trace level, this minimization is not a
single parity solve but a subset search with the parity decision as oracle
(Section~\ref{sec:complexity}); we state plainly that the strategy level delivers
validity via the game and minimality only via this separate, harder optimization.

\subsection{Worked liveness example}
\label{sec:liveness-example}

The toy effect \(F\,grant\) is reachability (co-safety) and needs no recurrence.
For a parity effect, take a two-state machine over input \(req\), output \(grant\):
\[
  \delta(\mathit{serve},\cdot)=\mathit{wait},\ \
  \delta(\mathit{wait},req)=\mathit{serve},\ \
  \delta(\mathit{wait},\neg req)=\mathit{wait};\quad
  \lambda(\mathit{serve},\cdot)=\{grant\},\ \lambda(\mathit{wait},\cdot)=\emptyset.
\]
Being Moore, the machine emits \(grant\) exactly \emph{while} in \(\mathit{serve}\)
(not the step after), entering \(\mathit{serve}\) one step after a request in
\(\mathit{wait}\). The initial state is \(s_0=\mathit{wait}\). Effect \(E=GF\,grant\); actual trace
\(\tau=(\{req\},\{grant\})^\omega\) (i.e.\ \(a_0=\{req\}\), \(a_1=\{grant\}\),
\dots), which grants infinitely often.

\paragraph{Fairness is required.} No controller guarantees \(GF\,grant\) against
\emph{every} environment: one issuing finitely many requests traps the machine in
\(\mathit{wait}\). The well-posed question carries the environment-fairness background
\(B_{\Frame}=GF\,req\), and the guaranteed objective is \(GF\,req\rightarrow
GF\,grant\), again \(\omega\)-regular. Without \(B_{\Frame}\) the winning region is
empty and every candidate is vacuously ``necessary,'' so stating \(B_{\Frame}\) is
what makes the example meaningful.

The monitor \(A_E\) for \(GF\,grant\) is a two-priority parity (Büchi) automaton
(priority \(0\) on \(grant\), \(1\) elsewhere; accepting iff \(0\) recurs); a
single priority cannot encode ``infinitely often.'' The product is accepting iff
the play visits \(\mathit{serve}\) infinitely often --- genuine recurrence, not
reachability.

\paragraph{The strategy-level cause is the recurrence commitment.} Take
\(C=GF\,\mathit{at\_serve}\) (effect-disjoint: it names \(\mathit{at\_serve}\), and
\(\Lang(GF\,\mathit{at\_serve})\not\subseteq\Lang(GF\,grant)\) by pure logic,
coinciding only through \(M\)). Under \(\Frame_{\mathrm{strat}}\) with
\(V_{\Frame}\) the \(\mathit{wait}\)-vertices and \(B_{\Frame}=GF\,req\), apply
Proposition~\ref{prop:rpg-sound}. The witness search asks for a play in
\(\mathcal{A}[W_\exists]\) satisfying \(GF\,req\wedge GF\,grant\wedge\neg C\),
i.e.\ \(GF\,req\wedge GF\,grant\wedge FG\,\mathit{at\_wait}\). But \(grant\) is
emitted only in \(\mathit{serve}\), so \(GF\,grant\) entails \(GF\,\mathit{at\_serve}\),
which contradicts \(FG\,\mathit{at\_wait}\); no such play exists, so \(C\) is
strategy-necessary. (Note this uses the \(B_{\Frame}\)-relativized N2: an unfair
environment with finitely many requests would, against the frozen \(\sigma\),
yield a play with \(FG\,\mathit{at\_wait}=\neg C\), but that play violates
\(B_{\Frame}=GF\,req\) and is therefore not an admissible deviation witness; this
is exactly why clause~(ii) is relativized to \(B_{\Frame}\).) The dominion \(D_C\)
is the \(\mathit{wait}\)-vertices from which \(\sigma\) serves on a request --- the
accepting-recurrence-preserving commitment. By contrast the \emph{stronger} ``serve
every request'' \(G((\mathit{at\_wait}\land req)\to X\,\mathit{at\_serve})\) is
\emph{not} necessary (a deviation may skip some requests yet serve infinitely
many, securing the objective on a \(B_{\Frame}\)-fair play while violating the
invariant), and the reachability candidate ``\(grant\) at time \(1\)'' is likewise
not necessary. Only a recurrence-preserving commitment passes --- why parity, not
reachability, is the right acceptance here.

\section{Overdetermination, Preemption, and Necessity vs.\ Sufficiency}
\label{sec:overdetermination}

Pure necessity has known limitations under overdetermination and preemption; we
address them head-on.

\paragraph{Overdetermination.} If inputs \(a\) and \(b\) each force \(E=X\,grant\)
and both hold, neither singleton is necessary (removing \(a\) leaves \(b\), which
still forces the grant), so neither is valid --- the correct report under
necessity. The \emph{necessary condition} (not necessarily ``the actual
cause'' in a stronger sense) is the disjunction \(C=a\lor b\), whose closest
removal sets both false and destroys \(E\) --- \emph{provided}
\(\mathcal{K}_{\Frame}\) admits disjunctions. This does not conflict with the
readability-first preference order (Section~\ref{sec:formal}), which penalizes
disjunctions on size: \(a\lor b\) is dispreferred to any smaller valid singleton,
so it is returned \emph{only} when no smaller valid cause exists --- which is
exactly the overdetermined case, where there genuinely is no smaller necessary
cause. If \(\mathcal{K}_{\Frame}\) is
restricted (e.g.\ to literals), no singleton is valid and the method reports
\(\mathsf{Causes}=\emptyset\); the remedy is to enlarge \(\mathcal{K}_{\Frame}\).
In this symmetric example \(C^{\flat}\) coincides with \(a\lor b\) on the
admissible set, but this identification is special to the example: in general
\(C^{\flat}\) is the distance-defined neighborhood and need not equal any readable
disjunction.

\paragraph{Preemption.} If \(a\) forces \(E\) and \(b\) would have only had \(a\)
been absent, then on a trace with both, \(a\) alone is not necessary; the
necessary condition is again \(a\lor b\). Distinguishing preempting from preempted
needs sufficiency/contingency, which necessity does not provide. A frame wanting
this imports a Halpern--Pearl-style contingency by pinning \(b\) in
\(B_{\Frame}\), recovering \(a\) relative to that frame.

\paragraph{Why necessity.} Two objections. ``Necessity flags trivial conditions
(`the system was powered on').'' It does not: non-vacuity discards conditions no
admissible counterfactual can violate, and \(B_{\Frame}\) absorbs standing
operating assumptions into the setting. ``Why not sufficiency (Halpern--Pearl)?''
HP adds a minimal-sufficient-set requirement with a contingency quantifier; even
in the Boolean setting its complexity is definition-sensitive --- NP-complete for
the original definition in binary models, \(\Sigma_2^P\)-complete for the original
definition in general models, \(\mathrm{D}_2^P\)-complete for the updated
definition (Eiter--Lukasiewicz; Aleksandrowicz--Chockler--Halpern--Ivrii). Lifting
any of these to \(\omega\)-regular effects layers the contingency quantifier atop
language operations, so HP sufficiency is at least as expensive as --- and
conceptually heavier than --- our necessity test (a one-player shortest-edit plus
non-emptiness check, or, at the strategy level, a parity solve followed by a
non-emptiness check). We claim only that necessity is \emph{no harder} (and uniform and
decidable), not that sufficiency is \emph{strictly} harder over \(\omega\)-regular
effects (which would require fixing an HP-over-\(\omega\)-regular semantics we do
not develop). We thus target necessity by design, report overdetermined effects as
disjunctive causes when expressible, and expose contingency through \(B_{\Frame}\).

\section{Decidability and Complexity}
\label{sec:complexity}

We give conditional bounds, not a fully proved complexity theorem; hypotheses are
explicit. Let \(|M|,|A_E|,|A_{\neg C}|\) be the machine/monitor sizes and \(d\) the
priority count.

\paragraph{One-player (input-only / output-intervention).} Candidate checking is
non-vacuity (non-emptiness of \(\Adm_M^{\Frame}(\tau)\cap\Lang(\neg C)\) over
finite-change removals) plus effect-failure. We do \emph{not} assume the closeness
\emph{relation} is automatic (its primary key compares unbounded counts). Instead,
by Remark~\ref{rem:lasso-close}, the closest removal is a shortest-edit path in the
finite product \(M\times A_{\neg C}\times A_E\) (size \(O(|M|\,|A_{\neg C}|\,
|A_E|)\)), computable by BFS, then \(E\) is evaluated on it --- polynomial in the
product (plus parity solving \(|V|^{O(d)}\), quasi-polynomial, if \(E\) is a
genuine parity effect). For LTL candidates the \emph{deterministic} parity
monitors used in the explicit product are worst-case \emph{doubly} exponential in
\(|C|\) (\(\LTL\to\) deterministic parity is doubly exponential), so the explicit
shortest-edit product route is doubly exponential in the formula size. The
underlying language checks (non-emptiness of \(\Adm\cap\Lang(\neg C)\) and
effect-failure) are themselves no harder than \(\LTL\) model checking ---
\textsc{PSpace} in the formula size --- when carried out via \emph{nondeterministic}
monitors rather than explicit determinization; we do \emph{not} claim a uniform
\textsc{PSpace} bound for the full minimum-count validity quantification.
\(C^{\flat}\) is \(\omega\)-regular and computable when
\(\uparrow_{\preceq_\tau}(F)\) is (via the shortest-edit construction on lassos;
else an explicit automaticity hypothesis).

\paragraph{Minimal-cause search.} Over a finite template family,
\(|\mathcal{K}_{\Frame}|\) checks; over the open LTL grammar, bounded-size search
is decidable but exponential in the bound (as in syntax-guided synthesis).

\paragraph{Two-player (strategy-structure).} Deciding N2 for a \emph{fixed}
\(V_{\Frame}\) (Proposition~\ref{prop:rpg-sound}) is: solve the parity game for
objective \(B_{\Frame}\rightarrow E\) (size \(O(|V|\,|A_B|\,|A_E|)\), \(O(d)\)
priorities) to get the winning region \(W_\exists\), then a one-player
\(\omega\)-regular non-emptiness check for \(B_{\Frame}\wedge E\wedge\neg C\) on
\(\mathcal{A}[W_\exists]\) (product with \(A_B,A_E,A_{\neg C}\); an accepting-lasso
search). Parity solving is in \(\mathrm{UP}\cap\mathrm{co\text{-}UP}\) and
quasi-polynomial, and the non-emptiness check is polynomial in the product, so the
decision is quasi-polynomial. (We do \emph{not} invoke a permissive-strategy
computation; the winning region plus the re-imposed objective suffices.) Finding a
\(\subseteq\)-minimal \(V_{\Frame}\) (the minimal dominion) is a \emph{separate}
search over \(\le 2^{|V_0|}\) vertex subsets, each tested by the above. We do not
prove a tight class, but record an upper bound: checking that a given
\(V_{\Frame}\) is necessity-inducing is one quasi-polynomial parity decision, and
checking minimality is a co-\(\mathrm{NP}\)-style universal quantification over
its proper subsets (each a parity decision); so the minimal-dominion decision sits
within a polynomial-hierarchy level relative to a parity oracle (e.g.\ a
\(\Sigma_2\)-type search), naturally encoded as SAT/SMT/MaxSAT over
vertex-selection variables with the parity decision as oracle. Exact completeness
is left open.

\paragraph{Finite-horizon (TempoBench).} \(\mathcal{K}_{\Frame_{\mathrm{TB}}}\) is
finite but \emph{exponential} (\(\le 2|I|(k+1)\) literals, hence
\(\le 2^{2|I|(k+1)}\) constraint sets). Checking one candidate is a finite
reachability check over \(A_{\mathrm{HOA}}\) on the window, polynomial in
\(|A_{\mathrm{HOA}}|\) and \(k\). Validity is not monotone in \(\Gamma\)
(Remark~\ref{rem:nonmonotone}), so greedy deletion is unsound; a subset-minimal
cause needs exhaustive/hitting-set/CEGIS search. Each check is polynomial; the
search is the hard part.

\section{TempoBench Specialization}
\label{sec:tempobench}

We specialize the framework to a TempoBench-style temporal causal evaluation
(TCE) task. \emph{Conditionally on the assumptions (F1)--(F2) below, which we have
not verified against the released artifacts}, the solver receives the same
information and solves the same finite-horizon problem, returning a
\emph{necessity}-based answer. Beyond (F1)--(F2) there is a second gap: TempoBench
labels TCE with CORP, a Halpern--Pearl \emph{actual-causation} tool, whereas we
compute counterfactual \emph{necessity}, and the two differ on overdetermined
effects (Section~\ref{sec:overdetermination}). So this is a \emph{proposed,
necessity-based} alignment offering an \emph{alternate} object, not a demonstrated
reproduction of CORP's key. TempoBench has two tasks: temporal trace evaluation (TTE), asking whether a
finite trace is accepted by a supplied automaton, and TCE, asking for the minimal
input constraints necessary for a specified output at a specified
timestep.\footnote{TempoBench: \url{https://arxiv.org/abs/2510.27544};
implementation \url{https://github.com/nik-hz/tempobench}.}

\paragraph{What the solver is given.} An automaton \(A_{\mathrm{HOA}}\) in HOA
format, the AP list, a finite trace \(\tau_{\le n}=\alpha_0\cdots\alpha_n\)
(\(\alpha_t\subseteq I\cup O\)), and one or more point effects written e.g.\
\(\texttt{XXX }g=X^3 g\). The reverse question
\(M+\varphi+\Frame+E+\tau\to C^\star\) specializes to
\(A_{\mathrm{HOA}}+\Frame_{\mathrm{TB}}+\tau_{\le n}+(E_{\Frame}=X^k o)\to C^\star\):
\emph{which provided input facts were necessary for output \(o\) at timestep
\(k\)?} The effect class is
\(\mathcal{E}_{TB}=\{X^k o\mid o\in O_{\Frame},\ 0\le k<|\tau_{\le n}|,\
o\in\alpha_k\}\) (positive selected-output occurrences at fixed timesteps); broad
effects (\(F\,grant\), \(G(req\to F\,grant)\), \(GF\,ready\),
\(\mathrm{Acc}_{\mathrm{parity}}\), internal-state reachability) are extensions
beyond TCE.

\begin{remark}[What \(A_{\mathrm{HOA}}\) is; faithfulness assumptions]
\label{rem:hoa}
TempoBench presents the system as a finite-state \emph{acceptor}
\(A_{\mathrm{HOA}}=(Q,2^{I\cup O},\delta,q_0,F)\): it reads valuation sequences
(inputs and outputs together) and accepts a finite trace when its run from \(q_0\)
ends in \(F\) (TempoBench's TTE semantics). It is an acceptor over \(I\cup O\),
not an I/O transducer; \(\Lang(A_{\mathrm{HOA}})\) plays the role of \(\Lang(M)\),
not of the effect monitor. For faithfulness we assume, \emph{conditionally}:
\textbf{(F1)} \(\Lang(A_{\mathrm{HOA}})\) is exactly the system's behavior set
(not a specification monitor \(A_\varphi\)); and \textbf{(F2)} the behavior set is
\emph{functional in the inputs} (input-deterministic), since the automata come
from reactive synthesis --- consistent with the acceptor framing, and the property
the operational reading uses to speak of ``the'' output at each timestep.
\textbf{We have not independently verified (F1)--(F2) against the released HOA
artifacts}; whether they are input-deterministic system models rather than
specification monitors is a checkable benchmark property. Accordingly the
alignment is stated \emph{conditionally on (F1)--(F2)}: if they hold, the reduction
below is exact; if the supplied automata are specification monitors, (F1) fails and
the alignment does not apply. Given (F1)--(F2), the point effect \(X^k o\) needs no
separate monitor (it is the membership check \(o\in\alpha_k\)); there is no product
to build, and the solver uses \(A_{\mathrm{HOA}}\) directly.
\end{remark}

\paragraph{The specialized frame.} \(\Frame_{\mathrm{TB}}\) is input-only
(\(A_{\mathrm{HOA}}\) fixed; outputs, transitions, priorities not intervened on),
with candidates the timestep-indexed input literals over the window
\(0,\ldots,k\): constraint sets \(\Gamma\subseteq\{0,\ldots,k\}\times
\mathrm{Lit}(I)\) (\(\mathrm{Lit}(I)=\{i,\neg i\}\)), denoting
\(C_\Gamma=\bigwedge_{(t,\ell)\in\Gamma}X^t\ell\). There are \(\le 2|I|(k+1)\)
literals, hence \(\le 2^{2|I|(k+1)}\) candidate sets: the space is finite but
\emph{exponential}, not linear. The closeness order is the count-only coarsening,
counted in changed \emph{input cells}: under input-only intervention with (F2),
outputs are determined by inputs, so a counterfactual's chosen changes are its
input-cell changes, and we count those rather than the full-valuation positions of
the general order (Section~\ref{sec:similarity}). So \(\Close\) is the full set of
minimum-count removals; minimality of causes is by set inclusion over \(\Gamma\).

\paragraph{Finite-horizon definitions.} Entirely in finite-trace terms (no
infinitary acceptance for \(A_{\mathrm{HOA}}\) is used). A finite sequence
\(\beta=\beta_0\cdots\beta_n\) is \emph{accepted} iff its run from \(q_0\) ends in
\(F\) (under (F2) the run is unique given the inputs); write
\(\Lang_{\le n}(A_{\mathrm{HOA}})\) for accepted length-\((n{+}1)\) sequences and
\(\Adm^{\le k}_{A,\Frame_{\mathrm{TB}}}=\Lang_{\le n}(A_{\mathrm{HOA}})\). For a
candidate \(\Gamma\), a removal is an admissible \(\beta\) violating \(C_\Gamma\),
and
\[
  \Close^{\le k}_{\tau_{\le n},A}(\neg C_\Gamma)=
  \min\nolimits_{\#\text{changes}}\{\beta\in\Adm^{\le k}_{A,\Frame_{\mathrm{TB}}}
  \mid\beta\not\models C_\Gamma\}.
\]
\[
\begin{aligned}
  \mathsf{Causes}^{\le n}_{A,\Frame_{\mathrm{TB}},k}(X^k o)=\{\Gamma\mid{}&
  \tau_{\le n}\models C_\Gamma,\ o\in\alpha_k,\
  \Close^{\le k}_{\tau_{\le n},A}(\neg C_\Gamma)\neq\emptyset,\\
  &\forall\beta\in\Close^{\le k}_{\tau_{\le n},A}(\neg C_\Gamma).\ o\notin\beta_k\}.
\end{aligned}
\]
A minimal cause \(C^\star=C_{\Gamma_E^\star}\) uses any subset-minimal valid
\(\Gamma_E^\star\) (several may exist).

\begin{remark}[Validity is not monotone; greedy deletion is unsound]
\label{rem:nonmonotone}
Membership is not monotone under \(\subseteq\) on \(\Gamma\): enlarging \(\Gamma\)
can make a candidate valid (the cheapest removal must flip more) or invalid (the
closest removal flips an irrelevant added literal while keeping the output). So
``delete literals while valid'' need not reach a subset-minimal cause; computing
\(\Gamma_E^\star\) needs genuine search (exhaustive, hitting-set/MaxSAT, or CEGIS).
Each validity check is polynomial; the search is hard. Greedy becomes sound only
under an extra monotonicity hypothesis, which we do not assume.
\end{remark}

\begin{proposition}[Finite specialization --- a definitional correspondence]
\label{prop:tb-reduction}
This is a \emph{definitional} restatement, not a deep theorem: the finite-horizon
objects above were \emph{defined} to mirror the general clauses, and the content is
that the mirroring is exact and needs no infinitary acceptance. We state it as a
proposition only to record that correspondence precisely. With
\(\Frame_{\mathrm{TB}}\) as the
input-only frame, \(B_{\Frame}=\True\), admissible set
\(\Lang_{\le n}(A_{\mathrm{HOA}})\), the change-counting order on the window, subset
minimality, and (F1)--(F2), and interpreting the general \(\mathsf{Causes}\) clauses
over finite traces with \(\LTLf\) semantics, for every \(\Gamma\):
\(\Gamma\in\mathsf{Causes}^{\le n}_{A,\Frame_{\mathrm{TB}},k}(X^k o)\) iff
\(C_\Gamma\) satisfies the four valid-cause clauses on \(\tau_{\le n}\) under
\(\Frame_{\mathrm{TB}}\).
\end{proposition}
\begin{proof}
Each of the four clauses refers only to positions \(0,\ldots,k\) (actuality is
\(\tau_{\le n}\models C_\Gamma\) and \(o\in\alpha_k\); non-vacuity is non-emptiness
of admissible \(\neg C_\Gamma\)-removals; effect-failure asks \(o\notin\beta_k\)
over closest removals), each determined by the length-\((k{+}1)\) prefix under
\(\LTLf\). These are verbatim the clauses of
\(\mathsf{Causes}^{\le n}_{A,\Frame_{\mathrm{TB}},k}\): admissible sets coincide,
the closeness order is in both cases the count-only coarsening of \(\preceq_\tau\)
on the window (so \(\Close\) is the full minimum-count set and effect-failure
quantifies over all of it), and (F2) makes ``\(o\) at \(k\)'' a well-defined
function of the admissible inputs. No infinite object appears, so no
\(\omega\)-acceptance is needed.
\end{proof}

\paragraph{Evaluation compatibility.} TempoBench scores at the AP level (partial
credit per atomic proposition at a timestep) and the TS level (a timestep counts
only if its whole constraint set matches). The output should therefore preserve
timestep structure: a flat formula \(X\,in_2\land X^3\neg in_0\) should also be
reported grouped (\(1\!:in_2\), \(3\!:\neg in_0\)). The pipeline should preserve
TempoBench's difficulty features (effect depth \(k\), automaton states, transition
count, trace length, causal-input count, unique inputs) and add no hidden
information from the original specification.

\section{Limitations and Conclusion}
\label{sec:conclusion}

\paragraph{Established.} A frame-relative, counterfactual-necessity notion of
temporal cause, decided in the common regimes by automaton emptiness/shortest-edit
search and in the strategy regime by the reverse parity game with a
sound-and-complete decision (Proposition~\ref{prop:rpg-sound});
Proposition~\ref{prop:weakest} characterizes the nearest-failure certificate; and
Proposition~\ref{prop:tb-reduction} shows the finite construction is the literal
specialization of the general definition, aligning it with TempoBench TCE.

\paragraph{Scoped / conjectural.} Complexity statements (Section~\ref{sec:complexity})
are at the level of construction sizes and standard solving bounds, not a fully
proved complexity theorem; they rely on automaticity of \(\preceq_\tau\), which we
establish for the canonical order on lassos and otherwise assume. The reverse game
decides \emph{strategy-guarantee} necessity (N2), a distinct type-level notion from
trace-existential N1; we do not claim one game decides N1, nor that one game
performs the edited-vertex minimization. The HP comparison is a contrast, not a
proved separation. The TempoBench alignment is conditional on (F1)--(F2)
(Remark~\ref{rem:hoa}), which we have not verified against the artifacts.

\paragraph{Limitations.} The method reports frame-relative necessary conditions,
not full actual causes: disjunctions for overdetermined effects, no preempting/%
preempted discrimination, and \(\mathsf{Causes}=\emptyset\) when the necessary cause
lies outside \(\mathcal{K}_{\Frame}\). Outputs are determinate only after the
eligible vocabulary is declared. Evidence here is two worked examples plus the
TempoBench alignment; a non-toy case study --- a strategy-level extraction on a
realistic arbiter, and an end-to-end TempoBench TCE instance with
\(\Gamma^\star\) compared to ground truth --- is left to future work and is the
main thing needed before empirical claims.

\paragraph{Conclusion.} Frame-relative counterfactual necessity over a frozen
synthesized strategy, with a parity-game lift for strategy-level causes, is a
useful organizing frame for temporal causal explanation. This document fixes the
definitions, the two regimes, and their soundness relationships; closing the
remaining gaps (a proved complexity theorem, the dominion-minimization procedure,
and an empirical evaluation) would turn it into a complete method.

\appendix

\section{Compact Notation}
\label{app:symbols}

Standard logical, temporal, set, and automaton notation is used as usual. We list
only the document-specific or overloaded objects.
\begin{description}[leftmargin=2.6cm,style=nextline]
  \item[\(\Frame\)] the causal frame
        \((B_{\Frame},E_{\Frame},S_{\Frame},O_{\Frame},\mathcal{K}_{\Frame},
        \mathcal{I}_{\Frame},\preceq,\triangleleft)\); explanatory scope, not a
        Kripke frame and not \(\varphi\).
  \item[\(E_{\Frame}\)] the frame-selected effect (\(=E\) if no scoping); a single
        \(\omega\)-regular trace property.
  \item[\(\mathcal{K}_{\Frame}\)] candidate-cause space (effect-disjoint,
        Assumption~\ref{ass:anticirc}).
  \item[\(\mathcal{I}_{\Frame}\)] intervention model fixing
        \(\Adm_M^{\Frame}(\tau)\) (input-only / output-intervention /
        strategy-structure). Distinct from the input set \(I\).
  \item[\(\Adm_M^{\Frame}(\tau)\)] admissible counterfactual traces under the
        frame.
  \item[\(\preceq_\tau\)] similarity order; canonically the count-first
        lexicographic \emph{total} order (Section~\ref{sec:similarity}), unique
        minimum \(\tau\). \(\uparrow_{\preceq_\tau}(X)\): traces \emph{farther}
        from \(\tau\) than some element of \(X\).
  \item[\(\Close_{\tau,M}^{\Frame}(\neg C)\)] \(\preceq_\tau\)-minimal admissible
        \(\neg C\)-removals (over finite-change removals).
  \item[\(\triangleleft\)] preference order on causes; canonically
        (size, depth, permissiveness), \(C'\triangleleft C\) means \(C'\) preferred.
  \item[\(\mathsf{Causes}_{M,\Frame,\tau}(E_{\Frame})\)] valid causes
        (Section~\ref{sec:formal}); \(C^\star\) a \(\triangleleft\)-minimal one.
  \item[\(C^{\flat}\)] the nearest-failure certificate
        (Section~\ref{sec:certificate}); determined only on
        \(\Adm_M^{\Frame}(\tau)\), generally not in \(\mathcal{K}_{\Frame}\).
  \item[N1 / N2] trace-existential (one-player) / strategy-guarantee (two-player)
        necessity (Section~\ref{sec:two-notions}).
  \item[\(\Sigma_{\Frame},V_{\Frame}\)] strategies agreeing with \(\sigma\) outside
        the editable vertex set \(V_{\Frame}\); \(D_C,T_C\) the causal dominion and
        its transitions.
  \item[{\(G^{\mathrm{rev}}_{M,\Frame,C,E}\), \(W_\exists\), \(\mathcal{A}[W_\exists]\)}]
        the reverse parity game (Definition~\ref{def:rpg}), its
        \(\exists\)-winning region, and the sub-arena restricted to it
        (Proposition~\ref{prop:rpg-sound}).
  \item[\(A_{\mathrm{HOA}}\), \(\Gamma\), \(C_\Gamma\)] TempoBench acceptor; an
        input-literal constraint set and its formula (Section~\ref{sec:tempobench}).
  \item[\(\mathit{at\_s}\)] atomic proposition exposing internal state \(s\)
        (prefix \(\mathit{at\_}\) to avoid clash with input names).
\end{description}

\end{document}
